\chapter{Conduire le projet}

Les méthodes dites agiles viennent du développement logiciel : des itérations courtes, un
résultat utilisable à la fin de chacune, et un plan qu'on revoit souvent. Elles
s'adaptent bien à un projet mathématique long, à condition de les ramener à quelques règles.

\section{Des lots de taille fixe}

\begin{regle}
On découpe le travail en lots de taille fixe, choisie pour que la qualité ne baisse pas
d'un lot à l'autre.
\end{regle}

Pour une transcription, vingt pages ; pour une preuve, un énoncé ; pour une relecture, une
séance de vingt minutes. Un lot fini est un lot commité, avec son message. La taille fixe
rend le travail mesurable : on sait combien de lots restent, et à quelle vitesse ils se
font.

\section{Une chose à la fois}

Les lots d'un même dossier peuvent avancer en parallèle ; les dossiers, eux, se traitent
l'un après l'autre, chacun fini et commité avant le suivant. Un projet où dix chantiers
sont ouverts à moitié ne livre rien.

\section{Un plan qu'on revoit}

Le plan vit dans les issues : une issue par chantier, avec sa liste de cases. Chaque
semaine, on relit les issues ouvertes, on ferme ce qui est réglé, on note ce qui attend une
ressource extérieure (un serveur en panne, une réponse, une autorisation) et l'on choisit
ce qui avance. Avant de proposer une tâche, on vérifie qu'elle n'est pas déjà faite.

\section{Attendre les autres}

Un projet de recherche attend souvent : un serveur d'archives indisponible, un spécialiste
qui doit répondre, une revue qui délibère. On l'écrit dans l'issue, avec la date, et l'on
avance sur autre chose. On ne relance pas un serveur ou une personne plus souvent que
nécessaire.

\section{Postmortems et design docs}

Deux formats viennent des équipes d'ingénierie. Le \emph{design document} décrit, avant de
construire, ce qu'on veut faire, les options écartées et pourquoi : une feuille de route de
formalisation en est un. Le \emph{postmortem} raconte, après un incident, ce qui s'est
passé, sans chercher de coupable, et finit par les règles qu'on en tire. Deux sessions qui
poussent sur la même branche, un déploiement cassé par un fichier ignoré : chaque incident
du projet a laissé une règle écrite, et plusieurs règles de ce livre viennent de là.

\section{Écrire à des spécialistes}

Une question bien posée à la bonne personne vaut des semaines de recherche. On écrit
court, on donne la question précise, ce qu'on a déjà cherché, et la page en cause ; on
relance une fois, après un délai raisonnable, puis on n'insiste plus.
